\textbf{Phần 1. Trắc nghiệm lựa chọn}
\setcounterpageref{ex}{0}

\begin{ex} %[4H2N08]
Góc nhọn là góc
\choice
{\True Bé hơn góc vuông}
{Lớn hơn góc vuông}
{Bằng góc vuông}
{Bằng góc bẹt}
\loigiai{
Góc nhọn bé hơn góc vuông ($< 90^\circ$).\\
\textbf{Đáp án của câu hỏi là A}
}
\end{ex}

\begin{ex} %[4H2N08]
Góc tù là góc
\choice
{Bé hơn góc vuông}
{Lớn hơn góc bẹt}
{\True Lớn hơn góc vuông và bé hơn góc bẹt}
{Bằng hai góc vuông}
\loigiai{
Góc tù lớn hơn góc vuông và bé hơn góc bẹt ($90^\circ < \text{Góc tù} < 180^\circ$).\\
\textbf{Đáp án của câu hỏi là C}
}
\end{ex}

\begin{ex} %[4H2H08]
Trong các góc sau, góc nào là góc nhọn?
\choice
{Góc $90^\circ$}
{\True Góc $45^\circ$}
{Góc $110^\circ$}
{Góc $180^\circ$}
\loigiai{
Góc $45^\circ < 90^\circ$ là góc nhọn.\\
\textbf{Đáp án của câu hỏi là B}
}
\end{ex}

\begin{ex} %[4H2H08]
Trong các góc sau, góc nào là góc tù?
\choice
{\True Góc $135^\circ$}
{Góc $30^\circ$}
{Góc $90^\circ$}
{Góc $180^\circ$}
\loigiai{
Góc $135^\circ$ là góc tù ($90^\circ < 135^\circ < 180^\circ$).\\
\textbf{Đáp án của câu hỏi là A}
}
\end{ex}

\begin{ex} %[4H2H08]
Hình chữ nhật có bao nhiêu góc vuông?
\choice
{$2$ góc vuông}
{\True $4$ góc vuông}
{$3$ góc vuông}
{$0$ góc vuông}
\loigiai{
Hình chữ nhật có 4 góc vuông.\\
\textbf{Đáp án của câu hỏi là B}
}
\end{ex}

\begin{ex} %[4H2H08]
Một hình tam giác có thể có nhiều nhất bao nhiêu góc vuông?
\choice
{$2$ góc}
{$3$ góc}
{$0$ góc}
{\True $1$ góc}
\loigiai{
Một hình tam giác chỉ có thể có nhiều nhất 1 góc vuông.\\
\textbf{Đáp án của câu hỏi là D}
}
\end{ex}

\begin{ex} %[4H2V08]
Hình tam giác đều có 3 góc bằng nhau. Mỗi góc của tam giác đều là
\choice
{\True Góc nhọn (mỗi góc bằng $60^\circ$)}
{Góc vuông}
{Góc tù}
{Góc bẹt}
\loigiai{
Mỗi góc của tam giác đều bằng $60^\circ$ nên là góc nhọn.\\
\textbf{Đáp án của câu hỏi là A}
}
\end{ex}

\begin{ex} %[4H2V08]
Trong một hình tứ giác có 2 góc vuông và 1 góc nhọn. Góc còn lại phải là
\choice
{Góc nhọn}
{Góc vuông}
{Góc bẹt}
{\True Góc tù}
\loigiai{
Tổng 4 góc tứ giác bằng $360^\circ$. Với 2 góc vuông ($180^\circ$) và 1 góc nhọn ($< 90^\circ$) thì góc còn lại phải $> 90^\circ$ (góc tù).\\
\textbf{Đáp án của câu hỏi là D}
}
\end{ex}

\begin{ex} %[4H2V08]
Mở một chiếc compa sao cho hai chân compa tạo thành một góc lớn hơn góc vuông nhưng chưa thẳng hàng. Góc đó là
\choice
{Góc nhọn}
{Góc bẹt}
{\True Góc tù}
{Góc vuông}
\loigiai{
Góc lớn hơn góc vuông và chưa thẳng hàng là góc tù.\\
\textbf{Đáp án của câu hỏi là C}
}
\end{ex}

\begin{ex} %[4H2V08]
Trong một hình thoi không phải là hình vuông, có
\choice
{4 góc vuông}
{4 góc nhọn}
{4 góc tù}
{\True 2 góc nhọn và 2 góc tù}
\loigiai{
Hình thoi không phải hình vuông có 2 góc nhọn đối nhau và 2 góc tù đối nhau.\\
\textbf{Đáp án của câu hỏi là D}
}
\end{ex}

\begin{ex} %[4H2V08]
Cho 3 tia chung gốc $OA, OB, OC$ sao cho góc $AOB = 40^\circ$ và góc $BOC = 50^\circ$ (tia $OB$ nằm giữa $OA$ và $OC$). Góc $AOC$ là
\choice
{Góc nhọn}
{Góc tù}
{\True Góc vuông ($90^\circ$)}
{Góc bẹt}
\loigiai{
Góc $AOC = 40^\circ + 50^\circ = 90^\circ$ là góc vuông.\\
\textbf{Đáp án của câu hỏi là C}
}
\end{ex}

\begin{ex} %[4H2C08]
Cho hình thang vuông $ABCD$ vuông tại $A$ và $D$. Góc tại đỉnh $C$ là góc tù $110^\circ$. Góc tại đỉnh $B$ là
\choice
{$60^\circ$}
{\True $70^\circ$ (góc nhọn)}
{$80^\circ$}
{$90^\circ$}
\loigiai{
Hai góc $B$ và $C$ bù nhau: Góc $B = 180^\circ - 110^\circ = 70^\circ$ (góc nhọn).\\
\textbf{Đáp án của câu hỏi là B}
}
\end{ex}

\textbf{Phần 2. Trắc nghiệm đúng sai}
\setcounterpageref{ex}{0}

\begin{ex} %[4H2H08]
Xét tính đúng sai về phân loại các góc:
\choiceTF
{\True Góc nhọn $<$ Góc vuông $<$ Góc tù $<$ Góc bẹt.}
{Góc nhọn lớn hơn góc vuông.}
{\True Một hình tam giác luôn có ít nhất 2 góc nhọn.}
{\True Góc bẹt có số đo bằng $2$ lần góc vuông.}
\loigiai{
\begin{itemize}
    \item Thứ tự độ lớn các góc: Đúng.
    \item Góc nhọn bé hơn góc vuông. Sai.
    \item Tam giác luôn có ít nhất 2 góc nhọn. Đúng.
    \item $180^\circ = 2 \times 90^\circ$. Đúng.
\end{itemize}
\begin{itemize}
    \item a) \textbf{Đúng}.
    \item b) \textbf{Sai}.
    \item c) \textbf{Đúng}.
    \item d) \textbf{Đúng}.
\end{itemize}
}
\end{ex}

\begin{ex} %[4H2V08]
Cho hình vẽ tam giác $ABC$ có góc $A = 90^\circ$, góc $B = 30^\circ$, góc $C = 60^\circ$. Xét tính đúng sai:
\choiceTF
{\True Tam giác $ABC$ là tam giác vuông tại đỉnh $A$.}
{\True Góc đỉnh $B$ là góc nhọn.}
{Góc đỉnh $C$ là góc tù.}
{\True Góc đỉnh $C$ là góc nhọn.}
\loigiai{
\begin{itemize}
    \item Góc $A = 90^\circ$ nên vuông tại $A$. Đúng.
    \item Góc $B = 30^\circ < 90^\circ$ là góc nhọn. Đúng.
    \item Góc $C = 60^\circ < 90^\circ$ là góc nhọn (không phải góc tù). Ý 3 Sai, Ý 4 Đúng.
\end{itemize}
\begin{itemize}
    \item a) \textbf{Đúng}.
    \item b) \textbf{Đúng}.
    \item c) \textbf{Sai}.
    \item d) \textbf{Đúng}.
\end{itemize}
}
\end{ex}

\textbf{Phần 3. Trả lời ngắn}
\setcounterpageref{ex}{0}

\begin{ex} %[4H2H08]
Hình vuông có tất cả bao nhiêu góc vuông?
\shortans{4}
\loigiai{
\textbf{Đáp số:} $4$\\
Hình vuông có 4 góc vuông.
}
\end{ex}

\begin{ex} %[4H2H08]
Một góc tù có số đo lớn hơn góc vuông $25^\circ$. Tính số đo của góc tù đó (chỉ ghi số).
\shortans{115}
\loigiai{
\textbf{Đáp số:} $115$\\
$90 + 25 = 115^\circ$.
}
\end{ex}

\begin{ex} %[4H2V08]
Một góc nhọn có số đo bằng một phần ba góc bẹt. Tính số đo góc nhọn đó (chỉ ghi số).
\shortans{60}
\loigiai{
\textbf{Đáp số:} $60$\\
$180 : 3 = 60^\circ$.
}
\end{ex}

\begin{ex} %[4H2V08]
Trong các góc: $35^\circ, 90^\circ, 120^\circ, 75^\circ, 150^\circ, 180^\circ$, có bao nhiêu góc nhọn?
\shortans{2}
\loigiai{
\textbf{Đáp số:} $2$\\
Có 2 góc nhọn ($35^\circ$ và $75^\circ$).
}
\end{ex}

\begin{ex} %[4H2V08]
Trong các góc: $35^\circ, 90^\circ, 120^\circ, 75^\circ, 150^\circ, 180^\circ$, có bao nhiêu góc tù?
\shortans{2}
\loigiai{
\textbf{Đáp số:} $2$\\
Có 2 góc tù ($120^\circ$ và $150^\circ$).
}
\end{ex}

\begin{ex} %[4H2C08]
Hai góc có tổng số đo bằng $180^\circ$ (tạo thành góc bẹt). Biết góc thứ nhất là góc tù và lớn hơn góc thứ hai là $40^\circ$. Tính số đo góc thứ hai (theo độ).
\shortans{70}
\loigiai{
\textbf{Đáp số:} $70$\\
Góc thứ hai: $(180 - 40) : 2 = 70^\circ$.
}
\end{ex}
